\section{Discussion, Limitations, and Research Integrity}
\label{sec:discussion}

\paragraph{Intervention-conditioned novelty is a hypothesis.}
A value-of-information objective becomes substantively interesting
only when source acquisition and model modification are
\emph{coupled in observed decisions}. It is possible that a fixed
teacher, an error-rule lookup table, or a standard contextual
optimizer matches \spear{} once every baseline receives equal
compute. In that case, a defensible conclusion is that
parameter-aware acquisition adds no measurable advantage in the
tested setting.

\paragraph{Diagnostic quality is not intervention utility.}
A correct teacher explanation may fail to translate into a
transferable update; an explanation that merely paraphrases
known trace state cannot create additional Shannon information
about hidden state. Improvements from such text should be
attributed to computational representation or inductive bias,
not new external knowledge.

\paragraph{Model changes can damage old skills.}
Low-rank updates and adapter isolation do not guarantee retention.
Rollback, versioned checkpoints, and old-skill audits are necessary.
Policies learned from a small candidate library may also overfit
to a single backbone, simulator, or error taxonomy. Updates from
a mismatched parameter basis should be rejected.

\paragraph{Counterfactual measurement is expensive.}
Paired post-update evaluation may dominate total cost; repeated
validation can create selection bias and spuriously optimistic
acceptance. The study must disclose cumulative cost, failed
attempts, multiplicity handling, and the number of adaptive
choices made on a development set. Test leakage can arise through
retrieved documentation, teacher answers, and persistent
experience archives, not just conventional gradient training.

\paragraph{Safety and deployment considerations.}
Autonomous parameter updates may amplify tool misuse or
unsafe behavior. The initial study should run in sandboxed,
resettable environments with restricted tools, reviewable update
diffs, and strict rollback. Real deployments require additional
security, privacy, and regression audits.

\section{Conclusion}
\label{sec:conclusion}
We have specified \spear{} as a \emph{proposed} framework for
evidence-seeking, parameter-level improvement of LLM agents.
Its central question is whether an agent can learn \emph{what
to investigate in order to change its future parameters well}.
The hypothesis is operationalized as parameter-intervention
value of information, a constrained evidence-to-update mapping,
and counterfactual validation of persistent model changes.
The defining experiments compare those mechanisms against
simple source rules, fixed sources, existing update generators,
and frozen-meta controls at matched resources.
Whether the method yields genuine recursive self-improvement
remains an empirical question.
